% Project 2 report. Every number comes from assets/numbers.tex or assets/tables/,
% generated by scripts/report_assets.py from the stored results.
% Build: latexmk -xelatex -cd projects/2/report/report.tex
\documentclass[11pt,a4paper]{article}
\usepackage{fontspec}
\usepackage[margin=2.3cm]{geometry}
\usepackage{amsmath}
\usepackage{booktabs}
\usepackage{graphicx}
\usepackage{float}
\usepackage{etoolbox}
\usepackage[backend=biber,style=numeric,sorting=none]{biblatex}
\usepackage[hidelinks]{hyperref}
\usepackage{xurl}
\addbibresource{refs.bib}
\input{assets/numbers.tex}
\graphicspath{{assets/figures/}}
\newcommand{\tab}[1]{\resizebox{\ifdim\width>\linewidth\linewidth\else\width\fi}{!}{\input{assets/tables/#1.tex}}}
\setlength{\parskip}{0.4em}
\setlength{\parindent}{0pt}

\title{Predicting GitHub Actions failures before the run\\[0.3em]
\large A from-scratch decision tree on self-collected CI data}
\author{ACISE -- Project 2}
\date{}

\begin{document}
\maketitle

\section{Problem}

A CI run costs time and machines, and a developer waits for its result. If the
outcome of a run could be predicted from what is known when the run is created,
likely failures could be run first and likely passes could be given lower priority.
This project predicts whether a GitHub Actions run will fail, using only
information available before the run starts: the change, its author and the
history of the workflow.

Two targets are used. \textbf{T1} is every run. \textbf{T2} keeps only the runs
whose previous run in the same lineage (the same pull-request branch or pushed
branch) passed, plus the first run of each lineage. T1 contains the runs that
merely continue a known failure, which the previous outcome already predicts
(the Parrot baseline). T2 removes them, so it asks for the harder and more useful
prediction: a failure after a pass. \textbf{T2 is the main target}; T1 is shown
for comparison.

The model is a CART decision tree written from scratch \cite{breiman1984cart}. A
tree is on the course list, needs no scaling, handles the imbalanced label with
class weights, and explains every prediction as a short path of threshold tests.
A random forest built from the same tree measures how much accuracy the
explanation costs. Related work predicts build outcomes from history and change
features \cite{hassan2017esem,chen2020buildfast,saidani2020ist} and uses such
predictions to skip builds \cite{jin2020icse}; those studies use Travis CI data,
whereas this project collects GitHub Actions data itself.

\section{Data collection}

\paragraph{Sources.} \NSources{} workflows of \NRepos{} open-source repositories
were collected, each over every branch and pull request unless restricted to one
branch (Table~\ref{tab:sources}). The built tables span \WindowStart{} to
\WindowEnd{}. The collection date $C$ is the end of the last collected day; the
frozen models are trained on runs created before $C$ (\Cutoff{}, 00:00 UTC), and
all later runs are new to them.

\paragraph{API.} For each workflow the GitHub REST API is queried in windows of at
most one month (halved down to one hour while a window holds 1\,000 runs or more,
the API's cap):
\url{GET /repos/<repo>/actions/workflows/<id>/runs?created=<window>},
then \url{/actions/runs/<id>/attempts/1} for re-run runs and
\url{/attempts/1/jobs} for failed first attempts. Every response is cached, so
the collection can be resumed and repeated. Each repository is also cloned in full,
including \texttt{refs/pull/*/head}; heads that are no longer referenced are
fetched by SHA.

\paragraph{Rows and features.} One row per run of a code-carrying event (push, pull
request, merge queue). Features come from the Actions history (failure rates over
time windows, flaky re-runs, CI load, queue times) and from git (churn, files and
directories touched, test and source files, author experience, file failure
history). History features of a run read only runs that \emph{completed} before it
was created, inside its own workflow, so no future information leaks in. Author
identities are replaced by salted SHA-256 hashes; the salt never leaves the raw
cache.

\paragraph{Labels.} The label is the conclusion of the \emph{first} attempt
(a passing re-run does not hide a failure). Runs that were cancelled or ended
otherwise are dropped. A failure whose failed jobs all match a per-workflow list
of infrastructure jobs (or that has no failed job) is dropped as well: it did not
come from the code. Of \RunsCollected{} collected runs, \RunsKept{} are kept and
\InfraDropped{} infrastructure failures are removed. T2 has \RowsTwo{} runs
(\FailPctTwo\,\% failed), T1 \RowsOne{} (\FailPctOne\,\% failed).

\begin{table}[H]
\centering\small
\tab{sources}
\caption{Sources. \emph{runs}: collected (all events); \emph{kept}: runs with a
pass or fail label; \emph{infra dropped}: infrastructure failures removed. Names
carry the workflow id when a repository has two workflows.}
\label{tab:sources}
\end{table}

\section{Approach and implementation}

\paragraph{Pooling.} The \NSources{} sources are pooled and ordered by creation
time, as if they were one diverse workflow. History features are computed per
source before pooling. The first 20\,\% of the pooled rows form the development
segment, used for exploration and feature selection.

\paragraph{Feature selection.} On the development segment only: near-constant and
mostly missing columns are removed, every feature is ranked by relevance statistics
(ANOVA F and Kendall $\tau$ for numeric features, $\chi^2$ for binary ones, mutual
information for all), the ranking is
repeated on 20 block subsamples, features chosen in at least 60\,\% of them are
kept, and of two features correlated above $\rho = 0.8$ only the more relevant
stays. For T2, a feature that falls out of the top 20 in three or more CV training
folds is then removed (drift filter). The frozen lists hold \NFeatTwo{} (T2) and
\NFeatOne{} (T1) features (Appendix, Table~\ref{tab:features}).

\paragraph{CART.} Binary splits $x \le t$ in raw units, Gini or entropy impurity,
optional balanced class weights, minimal cost-complexity pruning. The grid covers
depth 3--12, minimum leaf size 5--20, both criteria and both weightings. The
pruning $\alpha$ and the decision threshold on $P(\text{fail})$ are chosen on the
last 20\,\% of each training window, maximising MCC; one threshold is used for the
whole pool. scikit-learn's tree is used only in the tests, as a reference: on the
same inputs it gives the same predictions. The random forest bags 100 of our trees
on bootstrap samples, with $\sqrt{p}$ features per node.

\paragraph{Evaluation.} Expanding-window cross-validation with 5 folds after the
development segment: each fold trains on all earlier runs and tests on the next
block, so no model sees the future. The metric is the Matthews correlation
coefficient (MCC) of the failure class. Baselines: \emph{Parrot} (the lineage's
previous outcome), \emph{majority class}, and \emph{source\_rate}, which predicts
failure when the run's source has a training failure rate above the threshold. A
variant chooses one threshold per source. Methods are compared with paired
Wilcoxon tests and the Vargha--Delaney $A_{12}$ effect size on (fold, source)
cells with at least 10 test failures.

\paragraph{Frozen models.} One CART and one forest per target are trained on all
runs before $C$ with the same protocol and saved with skops, next to a JSON
sidecar that holds the features, hyperparameters, threshold, training medians (to
impute new runs), source list and the file's SHA-256. Loading checks the hash and
the allowed types. The frozen T2 CART has depth \DepthTwo{} and threshold \ThrTwo{}
(\NTrainTwo{} training and \NValTwo{} validation runs).

\section{Validation}

\subsection{Cross-validation}

\begin{table}[H]
\centering\small
\tab{cv}
\caption{Test MCC on the pooled rows, mean $\pm$ std over the 5 folds, and mean
recall of the failure class.}
\label{tab:cv}
\end{table}

On T2 the CART reaches an MCC of \CVCartTwo{} and the forest \CVForestTwo{},
against \CVSourceRateTwo{} for \emph{source\_rate}; Parrot and majority are 0,
since every T2 run follows a pass (Table~\ref{tab:cv}). On T1 the CART
(\CVCartOne{}) is level with Parrot (\CVParrotOne{}) and the forest is ahead
(\CVForestOne{}). Per-source thresholds do not help on the pooled rows: they even
out the flag rate across sources and give up the difference in base rates between
them. Table~\ref{tab:cells} gives the paired tests on the (fold, source) cells.

\begin{table}[H]
\centering\small
\tab{cells}
\caption{Paired Wilcoxon test and $A_{12}$ on (fold, source) MCC cells. Cells of
one fold share a model and a time window, so they are not independent.}
\label{tab:cells}
\end{table}

\begin{figure}[H]
\centering
\includegraphics[width=0.85\linewidth]{tree_cv_box_T2.png}
\caption{Test MCC per fold and method, T2.}
\end{figure}

\subsection{Transfer between repositories (E1--E3)}

E2 is the pooled model, scored on each source. E1 trains a CART on one source's
own runs and feature list; E3 trains the pooled model without the source's
repository and tests on that source, so the repository is unseen. All three use
the same test rows. Over \CellsEtwoEone{} cells, the pooled model is level with
the per-source models (E2 $-$ E1: \DiffEtwoEone{}, $p = \PEtwoEone{}$,
$A_{12} = \AEtwoEone{}$). Leaving the repository out costs accuracy (E3 $-$ E1:
\DiffEthreeEone{}, $p = \PEthreeEone{}$, $A_{12} = \AEthreeEone{}$), see
Table~\ref{tab:transfer}.

\begin{table}[H]
\centering\small
\tab{transfer}
\caption{T2 test MCC per source, mean over the 5 folds: pooled model (E2), the
source's own model (E1), and the pooled model trained without the source's
repository (E3).}
\label{tab:transfer}
\end{table}

\subsection{Re-test on runs after the collection date}

The frozen models were applied, unchanged, to every run created after $C$
(\RetestRunsTwo{} T2 and \RetestRunsOne{} T1 runs). Their features read the stored
history plus earlier new runs only. On T2 the CART scores \RetestCartTwo{} and the
forest \RetestForestTwo{}; on T1, \RetestCartOne{} and \RetestForestOne{}
(Table~\ref{tab:retest}). These are at or above the CV means, so the first days
after $C$ show no decay. Per source (Table~\ref{tab:retestsources}), the scores
spread widely and the small sources have few failures, so a single run moves
their MCC a lot.

\begin{table}[H]
\centering\small
\tab{retest}
\caption{Re-test of the frozen models on runs created after $C$, pooled.}
\label{tab:retest}
\end{table}

\begin{table}[H]
\centering\small
\tab{retest_sources}
\caption{T2 re-test per source.}
\label{tab:retestsources}
\end{table}

\begin{figure}[H]
\centering
\includegraphics[width=0.85\linewidth]{live_retest_vs_cv_T2.png}
\caption{CV folds (boxes) and the re-test after $C$ (points), T2, pooled.}
\end{figure}

\subsection{Explained predictions and a live run}

\texttt{predict-run} fetches one run, builds its row from its own source's history,
and prints the CART's decision path in raw feature units. The validation notebook
shows one true positive, true negative, false positive and false negative for each
of the sources with the most re-test failures, and predicts a run that is still in
progress: the actual use case, where the outcome is not yet known.

\subsection{Sandbox (E4)}

\ifdefstring{\SandboxDone}{yes}{%
The sandbox is a branch of this project's own repository with a small Python
package, real tests and one workflow. The models have never seen it. After
\SandboxWarmup{} unbroken warm-up changes, \SandboxScored{} scored changes were
pushed (\SandboxBreakPct\,\% broken), every change shape both broken and unbroken.
A broken change holds either a code defect (a mutation kept only if a test catches
it) or a CI break in the workflow file. Every run was predicted before it ended;
\SandboxDropped{} runs could not be scored. On the \SandboxRuns{} scored runs
(\SandboxFails{} failed) the T2 CART reaches an MCC of \SandboxMCC{} (Parrot:
\SandboxParrotMCC{}), see Table~\ref{tab:sandbox}.

\begin{table}[H]
\centering\small
\tab{sandbox_scores}\\[0.5em]
\tab{sandbox_breakage}
\caption{Sandbox runs: scores of the T2 CART, Parrot and majority, and the runs per
breakage kind.}
\label{tab:sandbox}
\end{table}
}{%
The sandbox experiment had not been run when these assets were generated.
}

\section{Limitations}

\begin{itemize}
\item The CART grid stops at depth 12; deeper trees were not explored.
\item Several T2 sources stay near an MCC of 0 (Table~\ref{tab:transfer}): the
  pre-run signal does not carry over to every workflow.
\item The (fold, source) cells of one fold share a model and a time window, so the
  paired tests treat dependent samples as independent.
\item The event features (push or pull request) differ in mix between sources, so
  they may partly stand in for the source rather than describe the change.
\item Failures caused by flaky tests or the CI environment cannot be predicted from
  pre-run data; the infrastructure filter removes only the ones it recognises.
\item In the sandbox the history is short, so most history features are imputed
  with the training medians and the prediction rests mainly on change features.
\end{itemize}

\section{Future work}

Per-repository fine-tuning on top of the pooled model; job-level prediction, so
that only the jobs likely to fail run first; features from the build logs of
earlier failures, as in BuildFast \cite{chen2020buildfast}; and a cost model that
weighs a missed failure against a delayed pass, to choose the threshold for a
concrete skipping or ordering policy.

\printbibliography

\appendix
\section{Frozen feature lists}
\begin{table}[H]
\centering\small
\tab{features}
\caption{Frozen feature lists, in relevance order.}
\label{tab:features}
\end{table}

\end{document}
